\documentclass[
%  published,
  draft
]{agujournal2025}

%%%%%%%%%%%%%%%%%%%%%%%%%%%%%%%%%%%%%%%%%%%%%%%%%%%%%%%%%%%%
% CONTRIBUTOR INSTRUCTIONS:
% Please add text only in the section you are assigned to.
% Keep first contributions concise: 1--3 paragraphs or bullet points.
% Use comments for unresolved questions.
% Add references as \cite{key} and add BibTeX entries to references.bib.
% Avoid broad mission-priority statements unless directly tied to the roadmap goal.
% Keep the focus on interdisciplinary gaps, capabilities, benchmarks, and recommendations.
%%%%%%%%%%%%%%%%%%%%%%%%%%%%%%%%%%%%%%%%%%%%%%%%%%%%%%%%%%%%

\begin{document}

\journalname{Perspectives of Earth and Space Scientists}

\title{From Solar Variability to Planetary Response: A Roadmap for Predictive Sun-to-Planet Coupling}

\authors{%
  Author One\affil{1},
  Author Two\affil{2},
  Author Three\affil{1,3}
}

\affiliation{1}{Author affiliation to be added.}
\affiliation{2}{Author affiliation to be added.}
\affiliation{3}{Author affiliation to be added.}

\authoraddr{Corresponding author name, institution, mailing address, and email address to be added.}

\authorrunninghead{Ignored.}
\titlerunninghead{Ignored.}

\keypoints%
{Solar influence on Earth and planetary environments requires integrated Sun-to-planet coupling across timescales.}
{Progress depends on benchmark events, 3D drivers, coupled models, and uncertainty-aware open data/modeling workflows.}
{This paper provides an implementation-oriented roadmap focused on gaps and community capabilities.}

\maketitle

\begin{abstract}
\textit{TODO Abstract: Summarize the focused, implementation-oriented community roadmap for predictive Sun-to-planet coupling. Identify interdisciplinary gaps, required tools and capabilities, benchmark datasets, and community recommendations. State that this paper does not duplicate the Decadal Survey or set broad mission priorities \cite{placeholder-roadmap,placeholder-decadal}.}

\begin{plainlanguagesummary}
\textit{TODO Plain Language Summary: Write a broad-access paragraph explaining that solar variability affects Earth and other planetary environments through long-term changes and transient events, but prediction remains difficult because observations, models, and communities are often disconnected across the Sun-to-planet chain.}
\end{plainlanguagesummary}
\end{abstract}

%%%%%%%%%%%%%%%%%%%%%%%%%%%%%%%%%%%%%%%%%%%%%%%%%%%%%%%%%%%%
% SECTION 1 OWNER MARKER: MOTIVATION AND SCOPE
%%%%%%%%%%%%%%%%%%%%%%%%%%%%%%%%%%%%%%%%%%%%%%%%%%%%%%%%%%%%
\section{Motivation and Scope}
% SECTION OWNER(S):
% MAIN MESSAGE:
% Explain why this white paper is needed and what it aims to accomplish.
% KEY POINTS TO INCLUDE:
% - Solar influence acts across timescales: solar minimum, solar cycle, storms, CMEs, ICMEs, SEPs, and solar-wind structures.
% - Planetary impacts depend on the full chain from solar variability to heliospheric propagation to environmental response.
% - Current work is fragmented across solar, heliospheric, geospace, planetary, modeling, and data communities.
% - This paper is not a Decadal Survey replacement; it is an implementation-oriented community roadmap.
% - The paper focuses on interdisciplinary gaps, shared capabilities, benchmark datasets, and recommendations.
% SUGGESTED OUTPUT:
% - One-page introduction.
% - Figure 1: schematic of the Sun-to-planet chain across timescales.
% LINKS TO RECOMMENDATIONS:
% - Connect this section to the full recommendation set in Section 5.

\noindent\textit{Draft text to be added by section owner.}

\begin{figure}[h]
\centering
\fbox{\begin{minipage}[c][2.0in][c]{0.85\textwidth}
\centering
Figure 1 drafting space: conceptual Sun-to-planet chain across timescales.
\end{minipage}}
\caption{Conceptual Sun-to-planet chain across timescales. Solar variability and eruptions propagate through the heliosphere and drive geospace and planetary responses. The schematic should show long-term forcing, transient events, heliospheric structure, domain coupling, and environmental impacts.}
\label{fig:sun-to-planet-chain}
\end{figure}

%%%%%%%%%%%%%%%%%%%%%%%%%%%%%%%%%%%%%%%%%%%%%%%%%%%%%%%%%%%%
% SECTION 2 OWNER MARKER: CENTRAL INTERDISCIPLINARY GAPS
%%%%%%%%%%%%%%%%%%%%%%%%%%%%%%%%%%%%%%%%%%%%%%%%%%%%%%%%%%%%
\section{Central Interdisciplinary Gaps}
% SECTION OWNER(S):
% MAIN MESSAGE:
% Identify where the Sun-to-planet chain currently breaks.
% KEY POINTS TO INCLUDE:
% - Magnetic-field structure from the corona to heliosphere is poorly observed and modeled.
% - Line-of-sight projection effects and sparse measurements limit 3D driver characterization.
% - L1 and single-point measurements do not fully capture the spatial/temporal structure of the driver.
% - Heliophysical driver information is not consistently transferred into geospace and planetary response models.
% - Background state, preconditioning, and recovery are often treated separately from transient events.
% - Benchmark datasets, uncertainty propagation, and common response metrics are insufficiently developed.
% SUGGESTED OUTPUT:
% - Gap table organized by interface:
%   Sun <-> Heliosphere
%   Heliosphere <-> Earth/Geospace
%   Earth/Geospace <-> Planetary Environments
%   Data/Modeling/AI <-> All Domains
% LINKS TO RECOMMENDATIONS:
% - Recommendation 2: 3D reconstruction.
% - Recommendation 3: domain coupling.
% - Recommendation 4: benchmark events and intervals.
% - Recommendation 5: validation and uncertainty.

\noindent\textit{Draft text to be added by section owner. Include placeholder citations where needed, for example \cite{placeholder-roadmap}.}

\begin{table}[h]
\caption{Interdisciplinary gaps in the Sun-to-planet chain.}
\label{tab:interdisciplinary-gaps}
\centering
\begin{tabular}{p{0.20\textwidth} p{0.26\textwidth} p{0.24\textwidth} p{0.20\textwidth}}
\hline
Domain interface & Gap & Limitation & Needed capability \\
\hline
Sun <-> Heliosphere & Coronal and heliospheric magnetic-field structure is poorly constrained & Limits prediction of CME/ICME orientation and Bz & Coronal magnetic-field constraints and 3D reconstruction \\
Heliosphere <-> Earth/Geospace & Single-point upstream measurements do not capture full driver structure & Limits geospace response prediction & Multipoint observations and uncertainty-aware propagation \\
Data/Models <-> All Domains & Benchmark datasets and shared metadata are incomplete & Limits validation and model comparison & Community benchmark events/intervals and standard metadata \\
\hline
\end{tabular}
\end{table}

%%%%%%%%%%%%%%%%%%%%%%%%%%%%%%%%%%%%%%%%%%%%%%%%%%%%%%%%%%%%
% SECTION 3 OWNER MARKER: PRIORITY SCIENCE GOALS
%%%%%%%%%%%%%%%%%%%%%%%%%%%%%%%%%%%%%%%%%%%%%%%%%%%%%%%%%%%%
\section{Priority Science Goals}
% SECTION OWNER(S):
% MAIN MESSAGE:
% Present the science questions that organize the white paper.
% KEY POINTS TO INCLUDE:
% - Each science goal should connect at least two domains.
% - Each science goal should arise from the identified gaps.
% - Each science goal should lead naturally to tools/capabilities and recommendations.
% SUGGESTED OUTPUT:
% - 3--4 science goals with short rationale paragraphs.
% LINKS TO RECOMMENDATIONS:
% - Map each subsection explicitly to one or more recommendations in Section 5.

\noindent\textit{Draft text to be added by section owner.}

%%%%%%%%%%%%%%%%%%%%%%%%%%%%%%%%%%%%%%%%%%%%%%%%%%%%%%%%%%%%
% SECTION 3.1 OWNER MARKER: 3D HELIOSPHERIC DRIVERS
%%%%%%%%%%%%%%%%%%%%%%%%%%%%%%%%%%%%%%%%%%%%%%%%%%%%%%%%%%%%
\subsection{Inferring the 3D Plasma and Magnetic Structure of Heliospheric Drivers}
\noindent\textbf{Core question:} How can we infer and validate the 3D plasma and magnetic structure of solar and heliospheric transients from sparse observations, and how does that structure evolve through the heliosphere?

% OWNER(S):
% MAIN MESSAGE:
% Prediction fails early if the heliospheric driver is not characterized in 3D.
% KEY POINTS TO INCLUDE:
% - Line-of-sight projection effects.
% - CME/ICME magnetic-field orientation and structure.
% - Sparse remote-sensing and in-situ observations.
% - Magnetic-field transport from photosphere/corona into heliosphere.
% - Multipoint observations and reconstruction methods.
% SUGGESTED FIGURE/TABLE CONTRIBUTION:
% - Optional schematic of projected observations -> 3D reconstruction -> planetary forcing.
% LINKS TO RECOMMENDATIONS:
% - Recommendation 1: coordinated observations.
% - Recommendation 2: 3D reconstruction.
% - Recommendation 3: heliospheric simulation coupling.

\noindent\textit{Draft text to be added by section owner.}

%%%%%%%%%%%%%%%%%%%%%%%%%%%%%%%%%%%%%%%%%%%%%%%%%%%%%%%%%%%%
% SECTION 3.2 OWNER MARKER: GEOSPACE RESPONSE
%%%%%%%%%%%%%%%%%%%%%%%%%%%%%%%%%%%%%%%%%%%%%%%%%%%%%%%%%%%%
\subsection{From Heliospheric Forcing to Coupled Geospace Response}
\noindent\textbf{Core question:} How do solar-wind and CME properties, together with the preconditioned geospace state, control energy partitioning, plasma transport, and response across Earth's magnetosphere-ionosphere-thermosphere system?

% OWNER(S):
% MAIN MESSAGE:
% Even if the upstream driver is known, prediction requires understanding coupled geospace response.
% KEY POINTS TO INCLUDE:
% - Magnetosphere-ionosphere-thermosphere coupling.
% - Energy partitioning during storms.
% - Radiation belts and plasmasphere response.
% - L1-to-Earth transformation of the driver.
% - Preconditioning and internal state.
% SUGGESTED FIGURE/TABLE CONTRIBUTION:
% - Response-chain diagram or schematic.
% LINKS TO RECOMMENDATIONS:
% - Recommendation 3: coupled response models.
% - Recommendation 4: benchmark events and reduced products.
% - Recommendation 5: uncertainty-aware modeling.

\begin{figure}[h]
\centering
\fbox{\begin{minipage}[c][2.0in][c]{0.85\textwidth}
\centering
Figure 2 drafting space: Solar wind / ICME driver $\rightarrow$ magnetospheric coupling $\rightarrow$ ionospheric electrodynamics $\rightarrow$ thermospheric heating $\rightarrow$ radiation belts / plasma transport $\rightarrow$ operational impacts.
\end{minipage}}
\caption{Heliophysical driver to geospace response chain: Solar wind / ICME driver, magnetospheric coupling, ionospheric electrodynamics, thermospheric heating, radiation belts / plasma transport, and operational impacts.}
\label{fig:geospace-response-chain}
\end{figure}

\noindent\textit{Draft text to be added by section owner.}

%%%%%%%%%%%%%%%%%%%%%%%%%%%%%%%%%%%%%%%%%%%%%%%%%%%%%%%%%%%%
% SECTION 3.3 OWNER MARKER: BACKGROUND STATE
%%%%%%%%%%%%%%%%%%%%%%%%%%%%%%%%%%%%%%%%%%%%%%%%%%%%%%%%%%%%
\subsection{Background State, Preconditioning, and Solar Timescales}
\noindent\textbf{Core question:} How do solar-cycle and solar-minimum conditions shape the background state of Earth's upper atmosphere, geospace, and planetary environments, and how does that background state modulate transient impacts?

% OWNER(S):
% MAIN MESSAGE:
% Transient responses cannot be understood without background state and system memory.
% KEY POINTS TO INCLUDE:
% - Solar minimum and solar-cycle forcing.
% - Background thermosphere, ionosphere, magnetosphere, radiation belts, and plasmasphere.
% - Similar events producing different impacts.
% - Recovery and system memory.
% - Avoid overbroad unsupported climate claims; frame around upper atmosphere, geospace, and planetary environments.
% SUGGESTED FIGURE/TABLE CONTRIBUTION:
% - Timescale table.
% LINKS TO RECOMMENDATIONS:
% - Recommendation 4: benchmark intervals and solar-minimum intervals.
% - Recommendation 5: validation and uncertainty.

\begin{table}[h]
\caption{Solar forcing, background state, and response across timescales.}
\label{tab:timescale-background-response}
\centering
\begin{tabular}{p{0.15\textwidth} p{0.20\textwidth} p{0.22\textwidth} p{0.18\textwidth} p{0.17\textwidth}}
\hline
Timescale & Solar forcing & Response domain & Key unknown & Needed observations/models \\
\hline
Solar minimum & Low EUV / altered solar wind & Thermosphere-ionosphere & Baseline state & Long-term monitoring \\
Solar cycle & Irradiance / solar-wind variability & Geospace climatology & Long-term modulation & Coordinated observations \\
Storms & CMEs / ICMEs / shocks & Coupled geospace response & Event-scale impacts & Benchmark events \\
Recovery & Post-storm conditions & Plasmasphere/radiation belts & System memory & Time-continuous measurements \\
\hline
\end{tabular}
\end{table}

\noindent\textit{Draft text to be added by section owner.}

%%%%%%%%%%%%%%%%%%%%%%%%%%%%%%%%%%%%%%%%%%%%%%%%%%%%%%%%%%%%
% SECTION 3.4 OWNER MARKER: COMPARATIVE PLANETARY RESPONSE
%%%%%%%%%%%%%%%%%%%%%%%%%%%%%%%%%%%%%%%%%%%%%%%%%%%%%%%%%%%%
\subsection{Universal and Planet-Specific Responses to Solar Forcing}
\noindent\textbf{Core question:} Which responses to solar forcing are universal across planetary environments, and which are controlled by planet-specific properties?

% OWNER(S):
% MAIN MESSAGE:
% Comparative planetary environments help identify which processes generalize beyond Earth.
% KEY POINTS TO INCLUDE:
% - Planetary magnetic field, atmosphere, rotation, heliocentric distance, and plasma sources.
% - Internally driven magnetospheres.
% - Need for upstream solar-wind/IMF/EUV context at planets.
% - Common response metrics.
% - Keep concise; avoid becoming a separate planetary habitability review.
% SUGGESTED FIGURE/TABLE CONTRIBUTION:
% - Comparative response table.
% LINKS TO RECOMMENDATIONS:
% - Recommendation 1: planetary observations.
% - Recommendation 6: comparative datasets and metrics.

\begin{table}[h]
\caption{Comparative planetary response framework.}
\label{tab:comparative-planetary-response}
\centering
\begin{tabular}{p{0.17\textwidth} p{0.20\textwidth} p{0.20\textwidth} p{0.18\textwidth} p{0.17\textwidth}}
\hline
Environment & External driver & Internal state & Response metric & Transferable lesson \\
\hline
Earth & To be added & To be added & To be added & To be added \\
Mars & To be added & To be added & To be added & To be added \\
Jupiter/Saturn & To be added & To be added & To be added & To be added \\
Venus & To be added & To be added & To be added & To be added \\
\hline
\end{tabular}
\end{table}

\noindent\textit{Draft text to be added by section owner.}

%%%%%%%%%%%%%%%%%%%%%%%%%%%%%%%%%%%%%%%%%%%%%%%%%%%%%%%%%%%%
% SECTION 4 OWNER MARKER: CAPABILITIES
%%%%%%%%%%%%%%%%%%%%%%%%%%%%%%%%%%%%%%%%%%%%%%%%%%%%%%%%%%%%
\section{Capabilities Needed for Predictive Sun-to-Planet Coupling}
% SECTION OWNER(S):
% MAIN MESSAGE:
% Summarize the tool groups and capabilities identified by the breakout.
% KEY POINTS TO INCLUDE:
% - Observations -> 3D reconstructions -> heliospheric simulations -> domain coupling -> response modeling -> validation.
% - Physics-based modeling and AI tools should support the full workflow.
% - Benchmarks, reduced event products, uncertainty, and metadata are required across the workflow.
% SUGGESTED OUTPUT:
% - Capability map table.
% - Optional workflow figure.
% LINKS TO RECOMMENDATIONS:
% - This section should provide the technical basis for all recommendations in Section 5.

\begin{figure}[h]
\centering
\fbox{\begin{minipage}[c][2.0in][c]{0.85\textwidth}
\centering
Figure 3 drafting space: Observations $\rightarrow$ 3D reconstruction $\rightarrow$ heliospheric simulations $\rightarrow$ domain coupling $\rightarrow$ geospace/planetary response models $\rightarrow$ benchmark validation.
\end{minipage}}
\caption{End-to-end Sun-to-planet workflow. Heliospheric and planetary observations feed 3D reconstruction, heliospheric simulations, domain coupling, geospace/planetary response models, and benchmark validation. Physics-based modeling, AI tools, uncertainty quantification, and open data support all stages.}
\label{fig:end-to-end-workflow}
\end{figure}

\noindent\textit{Draft text to be added by section owner. Include placeholder citations for benchmark/modeling concepts where needed, for example \cite{placeholder-benchmark}.}

%%%%%%%%%%%%%%%%%%%%%%%%%%%%%%%%%%%%%%%%%%%%%%%%%%%%%%%%%%%%
% SECTION 4.1 OWNER MARKER: OBSERVATIONS
%%%%%%%%%%%%%%%%%%%%%%%%%%%%%%%%%%%%%%%%%%%%%%%%%%%%%%%%%%%%
\subsection{Coordinated Heliospheric and Planetary Observations}
% OWNER(S):
% MAIN MESSAGE:
% Define the observing capabilities needed across solar, heliospheric, geospace, and planetary domains.
% KEY POINTS TO INCLUDE:
% - Multiwavelength imaging.
% - Multipoint in-situ measurements.
% - Upstream solar-wind, IMF, and EUV context.
% - Cross-calibrated datasets.
% - Planetary observations and multimission datasets.
% SUGGESTED FIGURE/TABLE CONTRIBUTION:
% - Observational requirements list.
% LINKS TO RECOMMENDATIONS:
% - Recommendation 1.

\noindent\textit{Draft text to be added by section owner.}

%%%%%%%%%%%%%%%%%%%%%%%%%%%%%%%%%%%%%%%%%%%%%%%%%%%%%%%%%%%%
% SECTION 4.2 OWNER MARKER: RECONSTRUCTIONS
%%%%%%%%%%%%%%%%%%%%%%%%%%%%%%%%%%%%%%%%%%%%%%%%%%%%%%%%%%%%
\subsection{3D Reconstructions}
% OWNER(S):
% MAIN MESSAGE:
% Define reconstruction capabilities for plasma and magnetic structure from sparse and projected data.
% KEY POINTS TO INCLUDE:
% - Reconstruct plasma and magnetic structure from sparse and projected data.
% - Combine remote sensing and in-situ data.
% - Estimate uncertainties.
% - Provide inputs to simulations and response models.
% SUGGESTED FIGURE/TABLE CONTRIBUTION:
% - Reconstruction capability requirements.
% LINKS TO RECOMMENDATIONS:
% - Recommendation 2.

\noindent\textit{Draft text to be added by section owner.}

%%%%%%%%%%%%%%%%%%%%%%%%%%%%%%%%%%%%%%%%%%%%%%%%%%%%%%%%%%%%
% SECTION 4.3 OWNER MARKER: HELIOSPHERIC SIMULATIONS
%%%%%%%%%%%%%%%%%%%%%%%%%%%%%%%%%%%%%%%%%%%%%%%%%%%%%%%%%%%%
\subsection{Heliospheric Simulations}
% OWNER(S):
% MAIN MESSAGE:
% Define how simulations should connect reconstructed structures to planetary forcing.
% KEY POINTS TO INCLUDE:
% - MHD simulations from Sun to 1 AU and beyond.
% - Data-driven magnetic-field models.
% - Boundary conditions and validation.
% - Simulations that connect reconstructed structures to planetary forcing.
% SUGGESTED FIGURE/TABLE CONTRIBUTION:
% - Simulation requirements and coupling needs.
% LINKS TO RECOMMENDATIONS:
% - Recommendations 2 and 3.

\noindent\textit{Draft text to be added by section owner.}

%%%%%%%%%%%%%%%%%%%%%%%%%%%%%%%%%%%%%%%%%%%%%%%%%%%%%%%%%%%%
% SECTION 4.4 OWNER MARKER: DOMAIN COUPLING
%%%%%%%%%%%%%%%%%%%%%%%%%%%%%%%%%%%%%%%%%%%%%%%%%%%%%%%%%%%%
\subsection{Domain Coupling}
% OWNER(S):
% MAIN MESSAGE:
% Define interfaces between heliospheric simulations and geospace/planetary response models.
% KEY POINTS TO INCLUDE:
% - Coupling heliospheric simulations to geospace and planetary response models.
% - Standard driver quantities and metadata.
% - Boundary assumptions.
% - Spatial and temporal coupling across domains.
% SUGGESTED FIGURE/TABLE CONTRIBUTION:
% - Interface requirements between domains.
% LINKS TO RECOMMENDATIONS:
% - Recommendation 3.

\noindent\textit{Draft text to be added by section owner.}

%%%%%%%%%%%%%%%%%%%%%%%%%%%%%%%%%%%%%%%%%%%%%%%%%%%%%%%%%%%%
% SECTION 4.5 OWNER MARKER: BENCHMARKS
%%%%%%%%%%%%%%%%%%%%%%%%%%%%%%%%%%%%%%%%%%%%%%%%%%%%%%%%%%%%
\subsection{Benchmark Datasets and Reduced Event Products}
% OWNER(S):
% MAIN MESSAGE:
% Define shared benchmark datasets, intervals, and reduced products needed for comparison and validation.
% KEY POINTS TO INCLUDE:
% - Benchmark storm events.
% - Quiet-time and solar-minimum intervals.
% - Solar-cycle phase comparisons.
% - Event lists of CMEs, ICMEs, flares, and structures impacting Earth/planets.
% - Reduced datasets and assumptions for model comparison.
% SUGGESTED FIGURE/TABLE CONTRIBUTION:
% - Benchmark/event product recommendations.
% LINKS TO RECOMMENDATIONS:
% - Recommendation 4.

\noindent\textit{Draft text to be added by section owner.}

%%%%%%%%%%%%%%%%%%%%%%%%%%%%%%%%%%%%%%%%%%%%%%%%%%%%%%%%%%%%
% SECTION 4.6 OWNER MARKER: MODELING AI UNCERTAINTY
%%%%%%%%%%%%%%%%%%%%%%%%%%%%%%%%%%%%%%%%%%%%%%%%%%%%%%%%%%%%
\subsection{Physics-Based Modeling, AI Tools, and Uncertainty}
% OWNER(S):
% MAIN MESSAGE:
% Define how physics-based modeling, AI/ML, and uncertainty quantification should support prediction.
% KEY POINTS TO INCLUDE:
% - Physics-based coupled models.
% - Neural operators, foundation models, and other AI/ML tools.
% - Data assimilation.
% - Event detection and classification.
% - Model-error correction and surrogate modeling.
% - Physical interpretability and validation.
% - Uncertainty propagation.
% SUGGESTED FIGURE/TABLE CONTRIBUTION:
% - AI/modeling requirements and validation principles.
% LINKS TO RECOMMENDATIONS:
% - Recommendation 5.

\noindent\textit{Draft text to be added by section owner.}

\begin{table}[h]
\caption{Capability map.}
\label{tab:capability-map}
\centering
\begin{tabular}{p{0.18\textwidth} p{0.20\textwidth} p{0.20\textwidth} p{0.20\textwidth} p{0.14\textwidth}}
\hline
Science goal & Needed observations & Needed models/tools & Needed data products & Validation metric \\
\hline
3D heliospheric drivers & To be added & To be added & To be added & To be added \\
Coupled geospace response & To be added & To be added & To be added & To be added \\
Background state and preconditioning & To be added & To be added & To be added & To be added \\
Comparative planetary response & To be added & To be added & To be added & To be added \\
\hline
\end{tabular}
\end{table}

%%%%%%%%%%%%%%%%%%%%%%%%%%%%%%%%%%%%%%%%%%%%%%%%%%%%%%%%%%%%
% SECTION 5 OWNER MARKER: COMMUNITY RECOMMENDATIONS
%%%%%%%%%%%%%%%%%%%%%%%%%%%%%%%%%%%%%%%%%%%%%%%%%%%%%%%%%%%%
\section{Community Recommendations}
% SECTION OWNER(S):
% MAIN MESSAGE:
% Convert the science goals and tools into actionable recommendations.
% KEY POINTS TO INCLUDE:
% - Recommendations should be action-oriented.
% - Each recommendation should identify a concrete product or community action.
% - Each recommendation should map to one or more science goals.
% SUGGESTED OUTPUT:
% - Recommendation table with near-term actions.

\begin{description}
\item[Recommendation 1] Develop coordinated heliospheric and planetary observing strategies. Include multiwavelength imaging, multipoint in-situ measurements, upstream solar-wind/IMF/EUV monitoring, planetary response observations, cross-calibration, and metadata.
\item[Recommendation 2] Build 3D reconstruction capabilities for heliospheric plasma and magnetic structure. Include CME/ICME density and magnetic-field structure, remote-sensing plus in-situ fusion, uncertainty estimates, and validation events.
\item[Recommendation 3] Couple heliospheric simulations to geospace and planetary response models. Include standard driver quantities, model interfaces, coupling assumptions, and response metrics.
\item[Recommendation 4] Create benchmark events, benchmark intervals, and reduced data products. Include storm events, quiet-time and solar-minimum intervals, solar-cycle phase comparisons, event lists, reduced datasets, and validation metrics.
\item[Recommendation 5] Advance physics-based and AI-enabled modeling with validation and uncertainty. Include physics-informed ML, neural operators/foundation models, data assimilation, event detection, uncertainty-aware outputs, and benchmark validation.
\item[Recommendation 6] Develop comparative planetary response datasets and metrics. Include multimission planetary datasets, upstream context at planets, common response metrics, comparative response tables, and future observation requirements.
\end{description}

\begin{table}[h]
\caption{Community recommendations.}
\label{tab:community-recommendations}
\centering
\begin{tabular}{p{0.20\textwidth} p{0.16\textwidth} p{0.20\textwidth} p{0.20\textwidth} p{0.16\textwidth}}
\hline
Recommendation & Addresses science goal & Concrete product & Near-term action & Possible owner community \\
\hline
Coordinated observing strategies & To be added & To be added & To be added & To be added \\
3D reconstruction capabilities & To be added & To be added & To be added & To be added \\
Coupled simulations and response models & To be added & To be added & To be added & To be added \\
Benchmark events and reduced products & To be added & To be added & To be added & To be added \\
Validated AI-enabled modeling & To be added & To be added & To be added & To be added \\
Comparative planetary response metrics & To be added & To be added & To be added & To be added \\
\hline
\end{tabular}
\end{table}

\noindent\textit{Draft text to be added by section owner.}

%%%%%%%%%%%%%%%%%%%%%%%%%%%%%%%%%%%%%%%%%%%%%%%%%%%%%%%%%%%%
% SECTION 6 OWNER MARKER: IMPLEMENTATION PATH
%%%%%%%%%%%%%%%%%%%%%%%%%%%%%%%%%%%%%%%%%%%%%%%%%%%%%%%%%%%%
\section{Implementation Path and Next Steps}
% SECTION OWNER(S):
% MAIN MESSAGE:
% Explain how the community can move from this roadmap to action.
% KEY POINTS TO INCLUDE:
% - Near-term: define benchmark events/intervals, event lists, and metadata.
% - Near-term: assemble section owners and writing team.
% - Mid-term: build cross-domain datasets and model interfaces.
% - Mid-term: compare models and develop validation metrics.
% - Longer-term: mature end-to-end Sun-to-planet workflows.
% - Include community review process and target venue.
% SUGGESTED OUTPUT:
% - Implementation timeline.
% - Contributor/section-owner plan.
% LINKS TO RECOMMENDATIONS:
% - Convert Section 5 into a staged implementation plan.

\begin{table}[h]
\caption{Implementation path.}
\label{tab:implementation-path}
\centering
\begin{tabular}{p{0.15\textwidth} p{0.23\textwidth} p{0.20\textwidth} p{0.20\textwidth} p{0.14\textwidth}}
\hline
Timeframe & Action & Product & Participants & Success metric \\
\hline
0--3 months & Assemble writing team and outline & Draft white paper & Section owners & Complete v0 draft \\
3--6 months & Define benchmark event/interval candidates & Community list & Observers/modelers & Shared candidate set \\
6--12 months & Develop data/model interface requirements & Interface document & Domain experts & Prototype workflow \\
1--3 years & Model intercomparison and benchmark validation & Community benchmark study & Cross-domain teams & Published benchmark results \\
\hline
\end{tabular}
\end{table}

\noindent\textit{Draft text to be added by section owner.}

%%%%%%%%%%%%%%%%%%%%%%%%%%%%%%%%%%%%%%%%%%%%%%%%%%%%%%%%%%%%
% CONCLUSION OWNER MARKER
%%%%%%%%%%%%%%%%%%%%%%%%%%%%%%%%%%%%%%%%%%%%%%%%%%%%%%%%%%%%
\section{Conclusions}
% SECTION OWNER(S):
% MAIN MESSAGE:
% Restate the central message and the need for an integrated, benchmarked, uncertainty-aware Sun-to-planet workflow.
% KEY POINTS TO INCLUDE:
% - The paper's main contribution is a focused roadmap.
% - The key step is moving from disconnected domain-specific studies to coupled workflows.
% - Observations, reconstructions, simulations, domain coupling, benchmarks, and AI/physics-based tools should work together.
% - Community coordination is required.
% LINKS TO RECOMMENDATIONS:
% - Briefly restate how the recommendations support the main goal.

\noindent\textit{Draft text to be added by section owner.}

%%%%%%%%%%%%%%%%%%%%%%%%%%%%%%%%%%%%%%%%%%%%%%%%%%%%%%%%%%%%
% BACK MATTER
%%%%%%%%%%%%%%%%%%%%%%%%%%%%%%%%%%%%%%%%%%%%%%%%%%%%%%%%%%%%
\section*{Acknowledgments}
\textit{TODO Acknowledgments: Acknowledge the NASA 5th Eddy Cross-Disciplinary Symposium and Theme 1 breakout contributors, if appropriate. Add funding and institutional support.}

\section*{Data Availability Statement}
\textit{TODO Data Availability Statement: No new data were generated for this perspective article. Add details for Miro/notes availability if the group decides to share them.}

\section*{Conflicts of Interest}
\textit{TODO Conflicts of Interest: Add conflict-of-interest statement before submission.}

% References
% Add real citations as the manuscript develops.
\bibliography{references}

\end{document}
